% presentation/slides/06_gradient_em.tex
\begin{frame}{Descente de Gradient Directe vs Approche EM}
    Pour maximiser $l_n(\theta)$, deux choix algorithmiques s'opposent :
    
    \vspace{0.4cm}
    \begin{columns}[t]
        \begin{column}{0.48\textwidth}
            \setbeamercolor{block title}{bg=DarkBeige,fg=Cream}
            \begin{block}{\textbf{Descente de Gradient Directe}}
                \small
                \begin{itemize}
                    \item Calcul direct du gradient de la log-vraisemblance marginale.
                    \item Fait intervenir les poids a posteriori $w_{i,z}(\theta)$.
                    \item Sensible au choix du pas (Règle d'Armijo nécessaire).
                    \item \alert{Nécessite des stratégies multi-départs} pour échapper aux extrema locaux.
                \end{itemize}
            \end{block}
        \end{column}
        
        \begin{column}{0.48\textwidth}
            \setbeamercolor{block title}{bg=SoftGold,fg=Cream}
            \begin{block}{\textbf{Algorithme EM}}
                \small
                \begin{itemize}
                    \item Maximise une borne inférieure via les \textbf{données complètes} $(Y, Z)$.
                    \item Propriété de \textbf{croissance monotone} de la vraisemblance à chaque étape.
                    \item Découple la complexité : l'étape d'inversion redevient quadratique et linéaire.
                \end{itemize}
            \end{block}
        \end{column}
    \end{columns}
\end{frame}